\chapter{Compiler Architecture}
\label{ch:compiler}

This chapter describes the architecture and implementation of the \haira{} compiler.

\section{Overview}

The \haira{} compiler (\code{hairac}) transforms source code into native executables. It is written in Rust and uses Cranelift for code generation.

\begin{lstlisting}[language={},numbers=none]
Source (.haira)
       |
       v
   +--------+
   | Lexer  |  --> Tokens
   +--------+
       |
       v
   +--------+
   | Parser |  --> AST
   +--------+
       |
       v
   +----------+
   | Resolver |  --> Resolved AST (imports, symbols)
   +----------+
       |
       v
   +-------------+
   | Type Checker|  --> Typed AST
   +-------------+
       |
       v
   +----------+
   | AI Phase |  --> AI blocks expanded (optional)
   +----------+
       |
       v
   +----------+
   | Lowering |  --> HIR (High-level IR)
   +----------+
       |
       v
   +----------+
   | Codegen  |  --> Cranelift IR --> Native binary
   +----------+
\end{lstlisting}

\section{Lexer}

The lexer (scanner) converts source text into a stream of tokens.

\subsection{Token Types}

\begin{lstlisting}[language={}]
enum TokenKind {
    // Literals
    IntLit(i64),
    FloatLit(f64),
    StringLit(String),
    BoolLit(bool),

    // Identifiers and keywords
    Ident(String),
    Keyword(Keyword),

    // Operators
    Plus, Minus, Star, Slash, Percent,
    Eq, EqEq, NotEq, Lt, Gt, LtEq, GtEq,
    And, Or, Not, AndAnd, OrOr, Bang,
    PlusEq, MinusEq, StarEq, SlashEq,
    Pipe, DotDot, DotDotEq, Arrow, FatArrow,
    LeftArrow, Question,

    // Delimiters
    LParen, RParen, LBracket, RBracket, LBrace, RBrace,
    Comma, Colon, Dot, Semicolon,

    // Special
    Newline, Eof, Error(String),
}
\end{lstlisting}

\subsection{Lexer Implementation}

\begin{lstlisting}[language={}]
struct Lexer<'a> {
    source: &'a str,
    chars: Peekable<CharIndices<'a>>,
    line: usize,
    column: usize,
}

impl<'a> Lexer<'a> {
    fn next_token(&mut self) -> Token {
        self.skip_whitespace();

        match self.peek_char() {
            None => Token::new(TokenKind::Eof, self.span()),
            Some(c) => match c {
                '0'..='9' => self.scan_number(),
                '"' => self.scan_string(),
                'a'..='z' | 'A'..='Z' | '_' => self.scan_identifier(),
                '+' => self.scan_operator(TokenKind::Plus, TokenKind::PlusEq),
                // ... other cases
            }
        }
    }
}
\end{lstlisting}

\section{Parser}

The parser constructs an Abstract Syntax Tree (AST) from the token stream.

\subsection{AST Nodes}

\begin{lstlisting}[language={}]
enum Expr {
    Literal(Literal),
    Ident(Ident),
    Binary(Box<Expr>, BinOp, Box<Expr>),
    Unary(UnaryOp, Box<Expr>),
    Call(Box<Expr>, Vec<Expr>),
    FieldAccess(Box<Expr>, Ident),
    Index(Box<Expr>, Box<Expr>),
    If(Box<Expr>, Block, Option<Block>),
    Match(Box<Expr>, Vec<MatchArm>),
    Closure(Vec<Param>, Option<Type>, Block),
    Array(Vec<Expr>),
    Map(Vec<(Expr, Expr)>),
    Struct(Ident, Vec<(Ident, Expr)>),
    Range(Box<Expr>, Box<Expr>, bool),
    Pipe(Box<Expr>, Box<Expr>),
}

enum Stmt {
    VarDecl(Vec<Ident>, Option<Type>, Vec<Expr>),
    Assignment(Vec<LValue>, AssignOp, Vec<Expr>),
    Expr(Expr),
    If(Expr, Block, Option<Block>),
    For(Option<Ident>, Ident, Expr, Block),
    While(Expr, Block),
    Match(Expr, Vec<MatchArm>),
    Return(Option<Vec<Expr>>),
    Break(Option<Ident>),
    Continue(Option<Ident>),
    Defer(Box<Stmt>),
    Spawn(SpawnTarget),
    Block(Block),
}
\end{lstlisting}

\subsection{Recursive Descent}

The parser uses recursive descent with Pratt parsing for expressions:

\begin{lstlisting}[language={}]
impl Parser<'_> {
    fn parse_expression(&mut self) -> Result<Expr, ParseError> {
        self.parse_expr_bp(0)
    }

    fn parse_expr_bp(&mut self, min_bp: u8) -> Result<Expr, ParseError> {
        let mut lhs = self.parse_prefix()?;

        loop {
            let op = match self.peek() {
                Some(Token { kind: TokenKind::Plus, .. }) => BinOp::Add,
                // ... other operators
                _ => break,
            };

            let (l_bp, r_bp) = op.binding_power();
            if l_bp < min_bp {
                break;
            }

            self.advance();
            let rhs = self.parse_expr_bp(r_bp)?;
            lhs = Expr::Binary(Box::new(lhs), op, Box::new(rhs));
        }

        Ok(lhs)
    }
}
\end{lstlisting}

\section{Resolver}

The resolver performs:
\begin{itemize}
    \item Import resolution
    \item Symbol binding (variables, functions, types)
    \item Scope analysis
    \item Forward declaration handling
\end{itemize}

\subsection{Symbol Table}

\begin{lstlisting}[language={}]
struct SymbolTable {
    scopes: Vec<Scope>,
    current: usize,
}

struct Scope {
    symbols: HashMap<String, Symbol>,
    parent: Option<usize>,
}

enum Symbol {
    Variable { ty: TypeId, mutable: bool },
    Function { sig: FunctionSig },
    Type { def: TypeDef },
    Module { exports: HashMap<String, Symbol> },
}
\end{lstlisting}

\subsection{Import Resolution}

\begin{lstlisting}[language={}]
fn resolve_import(&mut self, import: &Import) -> Result<(), ResolveError> {
    let module = match &import.path {
        // Standard library
        path if is_stdlib(path) => self.load_stdlib(path)?,

        // Project-local
        path if is_local(path) => self.load_local(path)?,

        // External package
        path => self.load_external(path)?,
    };

    match &import.kind {
        ImportKind::Simple => {
            self.bind_module(import.alias_or_name(), module)?;
        }
        ImportKind::Selective(names) => {
            for name in names {
                let symbol = module.export(name)?;
                self.bind_symbol(name, symbol)?;
            }
        }
        ImportKind::Glob => {
            for (name, symbol) in module.exports() {
                self.bind_symbol(name, symbol)?;
            }
        }
    }

    Ok(())
}
\end{lstlisting}

\section{Type Checker}

The type checker verifies type correctness and performs type inference.

\subsection{Type Representation}

\begin{lstlisting}[language={}]
enum Type {
    // Primitives
    Int, I8, I16, I32, I64,
    UInt, U8, U16, U32, U64,
    Float, F32, F64,
    Bool,
    String,

    // Composites
    Array(Box<Type>),
    Map(Box<Type>, Box<Type>),
    Tuple(Vec<Type>),
    Option(Box<Type>),

    // Named
    Struct(StructId),
    Enum(EnumId),

    // Functions
    Function(Vec<Type>, Box<Type>),

    // Generics
    TypeVar(TypeVarId),
    Generic(GenericId, Vec<Type>),

    // Special
    Never,
    Error,
}
\end{lstlisting}

\subsection{Type Inference}

\begin{lstlisting}[language={}]
fn infer_expr(&mut self, expr: &Expr) -> Result<Type, TypeError> {
    match expr {
        Expr::Literal(lit) => Ok(self.literal_type(lit)),

        Expr::Ident(name) => {
            self.lookup_type(name)
        }

        Expr::Binary(lhs, op, rhs) => {
            let lhs_ty = self.infer_expr(lhs)?;
            let rhs_ty = self.infer_expr(rhs)?;
            self.binary_result_type(op, &lhs_ty, &rhs_ty)
        }

        Expr::Call(func, args) => {
            let func_ty = self.infer_expr(func)?;
            match func_ty {
                Type::Function(params, ret) => {
                    self.check_args(&params, args)?;
                    Ok(*ret)
                }
                _ => Err(TypeError::NotCallable(func_ty)),
            }
        }

        // ... other cases
    }
}
\end{lstlisting}

\subsection{Unification}

\begin{lstlisting}[language={}]
fn unify(&mut self, a: &Type, b: &Type) -> Result<Type, TypeError> {
    match (a, b) {
        // Identical types
        (a, b) if a == b => Ok(a.clone()),

        // Type variable
        (Type::TypeVar(id), t) | (t, Type::TypeVar(id)) => {
            self.bind_type_var(*id, t.clone());
            Ok(t.clone())
        }

        // Option types
        (Type::Option(a), Type::Option(b)) => {
            Ok(Type::Option(Box::new(self.unify(a, b)?)))
        }

        // Arrays
        (Type::Array(a), Type::Array(b)) => {
            Ok(Type::Array(Box::new(self.unify(a, b)?)))
        }

        // Mismatch
        _ => Err(TypeError::Mismatch(a.clone(), b.clone())),
    }
}
\end{lstlisting}

\section{AI Phase}

The AI phase expands \keyword{ai} blocks into regular functions.

\subsection{AI Block Processing}

\begin{lstlisting}[language={}]
fn process_ai_block(&mut self, ai: &AiBlock) -> Result<FnDecl, AiError> {
    // 1. Generate cache key
    let key = self.cache_key(ai);

    // 2. Check cache
    if let Some(cached) = self.cache.get(&key) {
        return self.parse_cir(cached);
    }

    // 3. Prepare context
    let context = AiContext {
        signature: self.format_signature(ai),
        types: self.collect_types(ai),
        intent: ai.intent.clone(),
    };

    // 4. Invoke backend
    let cir = self.backend.generate(&context)?;

    // 5. Validate CIR
    self.validate_cir(&cir, ai)?;

    // 6. Cache result
    self.cache.set(&key, &cir);

    // 7. Convert to AST
    self.cir_to_ast(&cir)
}
\end{lstlisting}

\subsection{Backend Interface}

\begin{lstlisting}[language={}]
trait AiBackend {
    fn generate(&self, context: &AiContext) -> Result<Cir, AiError>;
    fn name(&self) -> &str;
    fn is_available(&self) -> bool;
}

struct LlamaCppBackend {
    model_path: PathBuf,
    context_size: usize,
    // ...
}

struct OllamaBackend {
    host: String,
    model: String,
    // ...
}

struct AnthropicBackend {
    api_key: String,
    model: String,
    // ...
}
\end{lstlisting}

\section{Lowering}

Lowering transforms the typed AST into HIR (High-level IR).

\subsection{HIR Structure}

\begin{lstlisting}[language={}]
struct HirModule {
    functions: Vec<HirFunction>,
    types: Vec<HirType>,
    constants: Vec<HirConst>,
}

struct HirFunction {
    name: Symbol,
    params: Vec<HirParam>,
    return_type: HirType,
    body: Vec<HirStmt>,
    locals: Vec<HirLocal>,
}

enum HirStmt {
    Assign(LocalId, HirExpr),
    Store(HirPlace, HirExpr),
    Call(Option<LocalId>, FnId, Vec<HirExpr>),
    If(HirExpr, BlockId, Option<BlockId>),
    Loop(BlockId),
    Break(Option<LoopId>),
    Continue(Option<LoopId>),
    Return(Option<HirExpr>),
}
\end{lstlisting}

\subsection{Desugaring}

\begin{lstlisting}[language={}]
// Pipe operator
// a | f | g  -->  g(f(a))
fn desugar_pipe(&mut self, pipe: &Pipe) -> HirExpr {
    self.build_call(
        pipe.rhs,
        vec![self.desugar_expr(&pipe.lhs)]
    )
}

// String interpolation
// "Hello, ${name}!"  -->  "Hello, " + to_string(name) + "!"
fn desugar_interpolation(&mut self, s: &InterpolatedString) -> HirExpr {
    let parts: Vec<HirExpr> = s.parts.iter().map(|part| {
        match part {
            Part::Literal(s) => HirExpr::StringLit(s.clone()),
            Part::Expr(e) => self.build_to_string(self.lower_expr(e)),
        }
    }).collect();

    self.build_string_concat(parts)
}

// For loop
// for x in items { ... }  -->  iterator pattern
fn desugar_for(&mut self, f: &ForLoop) -> Vec<HirStmt> {
    let iter = self.build_iter_call(&f.iterable);
    let loop_var = self.fresh_local();

    vec![
        HirStmt::Assign(loop_var, iter),
        HirStmt::Loop(self.build_block(vec![
            HirStmt::Assign(f.var, self.build_next_call(loop_var)),
            HirStmt::If(
                self.build_is_some(f.var),
                self.lower_block(&f.body),
                Some(self.build_break_block()),
            ),
        ])),
    ]
}
\end{lstlisting}

\section{Code Generation}

Code generation uses Cranelift to produce native code.

\subsection{Cranelift Integration}

\begin{lstlisting}[language={}]
fn codegen_function(&mut self, func: &HirFunction) -> Result<(), CodegenError> {
    let mut builder = FunctionBuilder::new(&mut self.ctx.func, &mut self.builder_ctx);

    // Create entry block
    let entry = builder.create_block();
    builder.append_block_params_for_function_params(entry);
    builder.switch_to_block(entry);

    // Generate parameters
    for (i, param) in func.params.iter().enumerate() {
        let val = builder.block_params(entry)[i];
        self.locals.insert(param.id, val);
    }

    // Generate body
    for stmt in &func.body {
        self.codegen_stmt(&mut builder, stmt)?;
    }

    builder.seal_all_blocks();
    builder.finalize();

    Ok(())
}

fn codegen_expr(&mut self, builder: &mut FunctionBuilder, expr: &HirExpr) -> Value {
    match expr {
        HirExpr::IntLit(n) => builder.ins().iconst(types::I64, *n),

        HirExpr::Binary(lhs, op, rhs) => {
            let l = self.codegen_expr(builder, lhs);
            let r = self.codegen_expr(builder, rhs);
            match op {
                BinOp::Add => builder.ins().iadd(l, r),
                BinOp::Sub => builder.ins().isub(l, r),
                BinOp::Mul => builder.ins().imul(l, r),
                BinOp::Div => builder.ins().sdiv(l, r),
                // ...
            }
        }

        HirExpr::Call(func, args) => {
            let args: Vec<Value> = args.iter()
                .map(|a| self.codegen_expr(builder, a))
                .collect();
            let call = builder.ins().call(func.ir_ref, &args);
            builder.inst_results(call)[0]
        }

        // ...
    }
}
\end{lstlisting}

\section{Project Structure}

\begin{lstlisting}[language={},numbers=none]
haira/
  Cargo.toml
  src/
    main.rs               # CLI entry point
    lib.rs                # Library entry point

  crates/
    haira-lexer/          # Lexical analysis
    haira-parser/         # Parsing
    haira-ast/            # AST definitions
    haira-resolver/       # Name resolution
    haira-types/          # Type system
    haira-typeck/         # Type checking
    haira-ai/             # AI integration
    haira-cir/            # CIR handling
    haira-hir/            # High-level IR
    haira-codegen/        # Code generation
    haira-driver/         # Compilation driver
    haira-stdlib/         # Standard library
\end{lstlisting}

\section{Build Process}

\begin{lstlisting}[language=bash,numbers=none]
# Full build
haira build main.haira

# Debug build
haira build main.haira --debug

# Release build
haira build main.haira --release

# Check without codegen
haira check main.haira

# Run directly
haira run main.haira

# REPL
haira repl
\end{lstlisting}

\section{Error Messages}

\begin{lstlisting}[language={},numbers=none]
error[E0001]: type mismatch
  --> src/main.haira:15:12
   |
15 |     return "hello"
   |            ^^^^^^^ expected `int`, found `string`
   |
   = note: expected type `int`
              found type `string`

error[E0015]: undefined variable
  --> src/main.haira:20:5
   |
20 |     println(undefined_var)
   |             ^^^^^^^^^^^^^ not found in this scope
   |
   = help: did you mean `undefined_value`?

error[E0042]: missing return statement
  --> src/main.haira:25:1
   |
25 | fn calculate(x: int) -> int {
   | ^^^^^^^^^^^^^^^^^^^^^^^^^^^ this function must return a value
   |
26 |     x * 2
   |     ----- help: consider adding a `return`: `return x * 2`
\end{lstlisting}
